\section{Contenerización Docker}
\label{sec:docker}

La contenerización de Voctomix~2.0 persigue un doble objetivo: garantizar la
reproducibilidad del sistema en cualquier equipo eliminando el problema de dependencias
y versiones de sistema operativo, e implantar un mecanismo de aislamiento de procesos
que mejore la resiliencia ante fallos en producción. El despliegue se articula en dos
ficheros: un \texttt{Dockerfile} que define la imagen base del sistema y un
\texttt{docker-compose.yml} que orquesta el conjunto de servicios.

El código original del repositorio (año~2016) contenía un \texttt{Dockerfile} basado
en Ubuntu~15.10 (Wily) con dependencias incompatibles con los repositorios actuales y
una imagen de base \texttt{c3voc/voctomix} que ya no existe en Docker Hub. Se eliminó
completamente ese material y se redactó una arquitectura moderna desde cero.

\subsection{Dockerfile}

La imagen se construye sobre \texttt{ubuntu:22.04~LTS}, garantizando soporte a largo
plazo y parches de seguridad. La instalación de dependencias se realiza en modo
desatendido (\texttt{DEBIAN\_FRONTEND=noninteractive}) en una sola capa para minimizar
el tamaño de la imagen:

\begin{itemize}
  \item GStreamer~1.0 con sus plugins \textit{base}, \textit{good}, \textit{bad} y
        \textit{ugly}.
  \item FFmpeg, para la inyección de fuentes de vídeo.
  \item \textit{Bindings} Python de GObject (\texttt{python3-gi}), la biblioteca Cairo
        para renderizado de overlays, y la librería AMQP (\texttt{pika}) para la
        integración con RabbitMQ.
  \item Utilidades de red (\texttt{netcat}, \texttt{iproute2}) necesarias para el
        \textit{healthcheck} y el autodescubrimiento de IP.
  \item Librerías matemáticas \texttt{scipy}~\cite{virtanen2020scipy} y \texttt{numpy}~\cite{harris2020numpy}, necesarias para el
        cálculo de curvas B-Spline en las transiciones animadas.
\end{itemize}

El directorio de trabajo del contenedor es \texttt{/opt/voctomix} y el punto de
entrada es \texttt{voctocore.py}.

Un \textit{healthcheck} basado en \texttt{nc -zw1 localhost 9999} permite a Docker
Compose determinar cuándo el núcleo está listo para aceptar conexiones antes de lanzar
los servicios dependientes, evitando condiciones de carrera durante el arranque.

\subsection{Docker Compose y patrón \textit{Shared Network Namespace}}
\label{sec:docker_compose}

La arquitectura de Voctomix exige que las fuentes de cámara (procesos FFmpeg
independientes) se conecten al servidor local mediante \texttt{localhost}. En Docker,
cada contenedor tiene por defecto su propio \texttt{localhost} aislado, lo que rompería
todas las conexiones TCP de entrada al núcleo.

En lugar de reescribir las direcciones IP en todo el código fuente, se adoptó el patrón
\textbf{\textit{Shared Network Namespace}}: los contenedores de las cámaras y el Notario
declaran \texttt{network\_mode: "service:voctocore"}, lo que fusiona sus interfaces de
red virtuales con la del contenedor del núcleo. Desde el punto de vista de Linux, todos
estos procesos comparten la misma tarjeta de red, por lo que \texttt{localhost} es el
mismo en todos ellos. Este patrón es habitual en arquitecturas de \textit{sidecar} y
permite mantener el código de los scripts de cámara sin ninguna modificación.

El fichero \texttt{docker-compose.yml} define los siguientes servicios organizados en
una red \textit{bridge} privada (\texttt{voctomix\_net}):

\begin{table}[H]
  \centering
  \caption{Servicios definidos en \texttt{docker-compose.yml}.}
  \label{tab:servicios_docker}
  \begin{tabular}{|l|l|c|l|}
    \hline
    \textbf{Servicio} & \textbf{Imagen base} & \textbf{Puertos expuestos} & \textbf{Función} \\
    \hline
    voctocore      & local/voctomix & 9999, 9998/udp, 11000--12000, 14000--14005, 15000 & Núcleo multimedia \\
    cam1--cam4     & local/voctomix & (namespace compartido) & Fuentes de prueba FFmpeg \\
    stream\_blanker& local/voctomix & 17000--17001, 18000 & Bucles pausa/offline \\
    rabbitmq       & rabbitmq:mgmt  & 5672, 15672 & Broker AMQP \\
    notario        & local/voctomix & 8080 & Telemetría REST/AMQP \\
    \hline
  \end{tabular}
\end{table}

\subsection{Resiliencia ante condiciones de carrera}

Durante las pruebas se comprobó que los contenedores de cámara arrancaban
milisegundos antes de que el núcleo abriese los puertos TCP de entrada,
provocando fallos de conexión silenciosos. Se resolvió con dos mecanismos
complementarios:

\begin{itemize}
  \item \textbf{\texttt{restart: always}:} delega en el demonio de Docker la
        responsabilidad de reintentar la conexión hasta que voctocore esté listo.
        Si un contenedor de cámara falla, Docker lo resucita automáticamente en
        segundos sin intervención del operador.
  \item \textbf{Esperas activas:} en los scripts de arranque nativos, los
        \texttt{sleep} fijos se sustituyeron por bucles que consultan
        \texttt{ss -lnt} para confirmar que el puerto 9999 está en escucha antes
        de lanzar los servicios dependientes.
\end{itemize}

\subsection{Aislamiento híbrido: GUI nativa, servicios en contenedor}

Tras evaluar las opciones de virtualización de interfaces gráficas (X11/Wayland
en contenedores), se decidió mantener la GUI (\texttt{voctogui}) ejecutándose
de forma \textbf{nativa} en el sistema anfitrión, conectándose al núcleo Docker
a través del puerto 9999 expuesto. Esta arquitectura híbrida evita los graves
problemas de latencia y aceleración por hardware (GPU/VAAPI) que surgen al
intentar virtualizar servidores gráficos en contenedores.

Durante el desarrollo se experimentó con la aceleración por hardware VAAPI para
la compresión de previsualizaciones. El sistema intentaba delegar la codificación
de vídeo a la GPU mediante el elemento \texttt{vaapipostproc}, que no estaba
disponible en el entorno de pruebas. Se resolvió desactivando la directiva
\texttt{vaapi=mpeg2} en la configuración y forzando el procesamiento por CPU.

\subsection{Inyección de recursos mediante volúmenes}

Para que el código Python de voctocore encuentre los logos e imágenes de fondo en
la ruta que espera sin modificar el código fuente, se declaró el volumen
\texttt{./images:/opt/images} en el \texttt{docker-compose.yml}. Este mecanismo
actúa como un \textit{bind mount} que hace visibles los recursos del sistema
anfitrión dentro del contenedor en la ruta estática esperada por GStreamer
(el protocolo \texttt{file://} del elemento \texttt{gdkpixbufoverlay} requiere
rutas absolutas según el RFC~8089).

% --- Figura sugerida ---
% \begin{figure}[H]
%   \centering
%   \includegraphics[width=0.9\textwidth]{figuras/docker_compose_detalle.png}
%   \caption{Diagrama completo de servicios Docker Compose y red interna voctomix\_net.}
%   \label{fig:docker_compose_detalle}
% \end{figure}

\subsection{Patrón \textit{Clean Exit}: liberación atómica de recursos}
\label{sec:clean_exit}

En sistemas basados en microservicios y contenedores, los procesos suelen dejar
puertos en estado \texttt{TIME\_WAIT} al cerrarse, impidiendo un reinicio inmediato.
Para garantizar que el sistema sea \textbf{idempotente} (que se pueda arrancar y
parar infinitas veces sin dejar recursos bloqueados), se implementó la gestión
avanzada de señales del sistema operativo:

\begin{itemize}
  \item \textbf{En los scripts Bash:} la instrucción \texttt{trap} intercepta la señal
        \texttt{SIGINT} (Ctrl+C) y ejecuta una rutina de limpieza que llama a
        \texttt{docker compose down}, deteniendo todos los contenedores y liberando
        los puertos de forma ordenada.
  \item \textbf{En el Notario (Python):} la opción \texttt{SO\_REUSEADDR}
        (\texttt{HTTPServer.allow\_reuse\_address = True}) ordena al kernel que
        reasigne el puerto 8080 de forma inmediata tras un cierre, eliminando el
        \texttt{Errno~98} que aparecía en reinicios rápidos durante el desarrollo.
\end{itemize}
